% Part: sets-functions-relations
% Chapter: arithmetization
% Section: cauchy

\documentclass[../../../include/open-logic-section]{subfiles}

\begin{document}

\olfileid[tr]{sfr}{arith}{cauchy}

\olsection{Ek: Cauchy Dizileri Olarak Gerçel Sayılar}

\olref[cuts]{sec}'te gerçel sayıları Dedekind kesimleri olarak kurduk. Bu
kısımda alternatif bir kuruluşu açıklıyoruz. Bu kuruluş, Cauchy'nin (bugün)
Cauchy dizisi dediğimiz şeye ilişkin tanımına dayanır; ancak bu tanımın gerçel
sayıları \emph{kurmak} için kullanılması başta Weierstrass, Heine, M\'{e}ray
ve Cantor olmak üzere başka on dokuzuncu yüzyıl yazarlarına aittir. (Güzel
bir tarihçe için bkz.\ \citeauthor{OConnorRobertson:RN}
\citeyear{OConnorRobertson:RN}.)

On dokuzuncu yüzyıla gelmeden önce Simon Stevin'i (1548--1620) ele almak
yararlı olur. Kısaca Stevin, her gerçel sayıyı ondalık açılımı aracılığıyla
düşünebileceğimizi fark etti. Dolayısıyla $\sqrt{2}$ gibi irrasyonel bir
sayının bile şöyle başlayan düzgün bir ondalık açılımı vardır:
\[
	1.41421356237\ldots
\]
Ondalık açılımları küme teorisi içinde modellemek çok kolaydır: onları
yalnızca $d \colon \Nat \to \{0,\ldots,9\}$ fonksiyonları olarak ele alırız; burada
$d(n)$, ilgilendiğimiz $n$'inci ondalık basamaktır. Sonra gerçel sayının ondalık
ayıracından önce gelen kısmını (burada yalnızca $1$) ele almak için küçük bir
düzeltme yapmamız gerekir. Örneğin $0.999\ldots = 1$ olmasını güvence altına
almak için bir düzeltmeye (bir denklik bağıntısına) daha ihtiyacımız vardır.
Ancak gerçel sayıları Stevin'in tarzında küme teorisi içinde bütünüyle titiz
biçimde kurmak zor değildir.

Odağımız Stevin değil. (Stevin hakkında daha fazlası için
bkz.\ \citealt{KatzKatz2012}.) Ancak bununla yakından bağlantılı bir düşünce
şudur. $\sqrt{2}$'nin ondalık açılımını doğrudan ele almak yerine, giderek
daha hassas açılımları kullanarak $\sqrt{2}$'ye giderek daha çok yaklaşan
rasyonel sayıların oluşturduğu bir \emph{diziyi} ele alabiliriz:
\[
1, 1.4, 1.414, 1.4142, 1.41421,\ldots
\]
Gerçel sayıların “giderek daha iyi yaklaşımlar” aracılığıyla ele alınabileceği
fikri, bize başka bir gözlemler dizisinin temelini verir (bunlar $\Rat$'yi
$\Int$'den veya $\Int$'yi $\Nat$'den kurarken kullandığımız gözlemlere
benzer). Temel gözlemler şunlardır:
\begin{enumerate}
	\item Her gerçel sayı (belki sonsuz) bir ondalık açılım olarak yazılabilir.
	\item (Belki sonsuz) bir ondalık açılımda kodlanan bilgi, bir rasyonel
	sayılar dizisiyle de kodlanabilir.
	\item Bir rasyonel sayılar dizisi $\Nat$'den $\Rat$'ye bir fonksiyon olarak
	düşünülebilir; $f(n)$'yi dizideki $n$'inci rasyonel sayı olarak almak
	yeterlidir.
\end{enumerate}
Elbette $\Nat$'den $\Rat$'ye giden \emph{herhangi} bir fonksiyon bize bir gerçel
sayı vermez. Örneğin şu fonksiyonu ele alın:
\[
	f(n) =\begin{cases}
			1 & \text{eğer }n\text{ tekse}\\
			0 &\text{eğer }n\text{ çiftse}
		\end{cases}
\]
Buradaki temel kaygı, $0,1,0,1,0,1,0,\ldots$ dizisinin herhangi bir gerçel
sayıya “giderek yaklaşmıyor” görünmesidir. Dolayısıyla bir gerçel sayıya
yakınsayan dizileri ele aldığımızdan emin olmak için dikkatimizi bir
\emph{limiti} olan dizilerle sınırlandırmalıyız.

Limit fikriyle tarih bölümündeki “Limitlerin Kesin Tanımı” başlıklı kısımda
zaten karşılaştık. Ancak orada
kullandığımız tanımın \emph{tam olarak} aynısını kullanamayız. Oradaki
“$(\forall \epsilon>0)$” ifadesi, örtük olarak gerçel sayılar üzerinde
niceleme içeriyordu; ayrıca gerçel sayılar üzerindeki fonksiyonların limitlerini
ele alıyorduk. Bu nedenle o tanıma başvurmak, tam da \emph{kurmayı}
amaçladığımız gerçel sayıları hazır kabul etmek olurdu. Neyse ki yakından
bağlantılı bir limit fikriyle çalışabiliriz.
\begin{defn}\ollabel{def:CauchySequence}
  Bir $f: \Nat \to \Rat$ fonksiyonu, ancak ve ancak her pozitif
  $\epsilon \in \Rat$ için
  $(\exists \ell \in \Nat)(\forall m, n > \ell)|f(m) - f(n)| < \epsilon$
  ise bir \emph{Cauchy dizisidir}.
\end{defn}

Genel limit fikri öncekiyle aynıdır: belirli bir hassasiyet düzeyi
($\epsilon$ ile ölçülen) istediğinizde bakacağınız bir “bölge” (yani
$\ell$'den büyük her girdi) vardır. Ayrıca $1$, $1.4$, $1.414$, $1.4142$,
$1.41421$\ldots\ dizisinin bir limiti olduğunu görmek kolaydır:
$\sqrt{2}$'ye $\nicefrac{1}{10^{n}}$'den küçük bir hata payıyla yaklaşmak
istiyorsanız $n$'inci terimden sonraki herhangi bir terime bakmanız yeterlidir.

Öyleyse açık düşünce, bir gerçel sayının herhangi bir Cauchy dizisi
\emph{olduğunu} söylemek olurdu. Fakat $\Int$ ve $\Rat$ kuruluşlarında olduğu
gibi bu da fazla naif olurdu: verili herhangi bir gerçel sayıyı birden çok
farklı Cauchy dizisi gösterir. Bunu görmenin basit bir yolu şudur. Bir Cauchy
dizisi $f$ verildiğinde, $g(0)\neq f(0)$ olması dışında $f$ ile bütünüyle aynı
olan $g$ fonksiyonunu tanımlayın. İki dizi ilk sayıdan sonra her yerde
örtüştüğü için, sonunda \olref{def:CauchySequence}'de kullanılan anlamda aynı
limite sahip olduklarını söylemek isteyeceğiz; dolayısıyla aynı gerçel sayıyı
“tanımlıyor” kabul edilmelidirler. Öyleyse bu Cauchy dizilerini aynı gerçel sayı
olarak düşünmeliyiz.

Sonuç olarak Cauchy dizileri üzerinde yeniden bir denklik bağıntısı
tanımlamamız ve gerçel sayıları denklik sınıflarıyla özdeşleştirmemiz gerekir.
Önce limitte $0$'a yaklaşan fonksiyon fikrine ihtiyacımız var. Herhangi bir
$h : \Nat \to \Rat$ fonksiyonu için, ancak ve ancak her pozitif
$\epsilon \in \Rat$ için
$(\exists \ell \in \Nat)(\forall n > \ell)|h(n)| < \epsilon$ ise
\emph{$h$, $0$'a yakınsar} diyelim.\footnote{Bunu, tarih bölümündeki
“Limitlerin Kesin Tanımı” başlıklı kısımda verilen
$\lim_{x \mathord{\rightarrow}\infty}f(x) = 0$ tanımıyla karşılaştırın.}
Ayrıca $f$ ve $g$, $\Nat \to \Rat$ fonksiyonları olduğunda
$(f-g)(n) = f(n) - g(n)$ olsun. Şimdi şunu tanımlayalım:
\[
	f \Realequiv g \text{ ancak ve ancak $(f-g)$, $0$'a yakınsar.}
\]
$\Realequiv$ bağıntısının bir denklik bağıntısı olduğunu denetlemeliyiz; öyle
olduğu gösterilebilir. Ardından istersek gerçel sayıları, $\Nat \to \Rat$
biçimindeki bütün Cauchy dizilerinin $\Realequiv$ altındaki denklik sınıfları olarak
tanımlayabiliriz.

\begin{prob}
Her $n$ için $f(n) = 0$ olsun. $g(n) = \frac{1}{(n+1)^2}$ olsun. Her ikisinin
de Cauchy dizisi olduğunu ve gerçekten her iki fonksiyonun limitinin $0$
olduğunu, dolayısıyla $f \Realequiv g$ olduğunu gösterin.
\end{prob}

Bunu yaptıktan sonra, her zamanki gibi $f$'yi !!a{element} olarak içeren
denklik sınıfı için $\equivrep{f}{\Realequiv}$ yazacağız. Ancak okunabilirliği
korumak için bundan sonra alt indisi atıp yalnızca $\equivrep{f}{}$ yazacağız.
Ayrıca her $q \in \Rat$ için
$q_{\Real} = \equivrep{c_{q}}{}$ olduğunu belirliyoruz; burada $c_q$, bütün
$n \in \Nat$ için $c_q(n) = q$ olan sabit fonksiyondur. Sonra gerçel sayılar
üzerindeki temel bağıntı ve işlemleri tanımlarız; örneğin:
\begin{align*}
	\equivrep{f}{} + \equivrep{g}{} &= 	\equivrep{(f + g)}{} \\
	\equivrep{f}{} \times 	\equivrep{g}{} &= \equivrep{(f \times g)}{}
\end{align*}
burada $(f + g)(n) = f(n) + g(n)$ ve
$(f \times g)(n) = f(n) \times g(n)$'dir. Elbette $f$ ile $g$ Cauchy
dizileri olduğunda $(f + g)$, $(f-g)$ ve $(f\times g)$'nin her birinin de
Cauchy dizisi olduğunu denetlemeliyiz; gerçekten öyledirler ve bunu size
bırakıyoruz.

Son olarak bir sıralama kavramı tanımlarız. $\equivrep{f}{}$'nin, ancak ve
ancak hem $\equivrep{f}{}\neq 0_\Real$ hem de
$(\exists \ell \in \Nat)(\forall n > \ell)0 < f(n)$ ise \emph{pozitif}
olduğunu söyleyelim. Ardından $\equivrep{f}{} < \equivrep{g}{}$ bağıntısını,
ancak ve ancak $\equivrep{(g - f)}{}$ pozitifse geçerli sayalım. Bunun iyi
tanımlı olduğunu (yani denklik sınıfından seçilen “temsilci” fonksiyona bağlı
olmadığını) denetlemeliyiz.
%\begin{align*}
%	\equivrep{f}{} \leq \equivrep{g}{} \text{ iff }&\equivrep{f}{} = \equivrep{g}{} \text{ or }\\
%	&(\exists \ell \in \Nat)(\forall n \in \Nat)(\ell < n \lif f(n) \leq g(n))
%\end{align*}
Bunu yaptıktan sonra, bu tanımların doğru cebirsel özellikleri verdiğini
göstermek oldukça kolaydır; yani:
\begin{thm}\ollabel{thm:cauchyorderedfield}
Cauchy dizilerinin denklik sınıfları sıralı bir cisim oluşturur.
\end{thm}

\begin{proof}
Alıştırma.
\end{proof}

\begin{prob}
Cauchy dizilerinin denklik sınıflarının sıralı bir cisim oluşturduğunu
ispatlayın.
\end{prob}

Bu biçimde kurulan gerçel sayıların Tamlık Özelliği'ne sahip olduğunu ispatlamak
daha zordur; bu yüzden ispatı vereceğiz.

\begin{thm}
Boş olmayan bir Cauchy dizileri kümesinin karşılık gelen denklik sınıfları
üstten sınırlıysa, bu sınıfların en küçük bir üst sınırı vardır.
\end{thm}

\begin{proof}[İspat taslağı]
$S$, karşılık gelen denklik sınıfları üstten sınırlı olan, boş olmayan
herhangi bir Cauchy dizileri kümesi olsun. Öyleyse $p_{\Real}$, $S$'nin
denklik sınıfları için bir üst sınır olacak biçimde bir $p \in \Rat$ vardır.
$r \in S$ olsun; o zaman $q_{\Real} < \equivrep{r}{}$ olacak biçimde bir
$q \in \Rat$ vardır. Dolayısıyla $S$ üzerinde en küçük bir üst sınır varsa
bu sınır $q_\Real$ ile $p_\Real$ arasındadır (uçlar dâhil).

En küçük üst sınıra hem aşağıdan hem yukarıdan aynı anda yaklaşarak onu
belirleyeceğiz. Özellikle $f, g \colon \Nat \to \Rat$ biçiminde iki fonksiyon
tanımlarız; amacımız $f$'nin en küçük üst sınıra yukarıdan, $g$'nin ise
aşağıdan giderek yaklaşmasıdır. Şöyle başlarız:
\begin{align*}
	f(0) &= p \\
	g(0) &= q
\end{align*}
Sonra $a_n = \frac{f(n) + g(n)}{2}$ olmak üzere şunları
tanımlayalım:\footnote{Bu özyinelemeli bir tanımdır. Ancak özyinelemeli
tanımların meşru olduğunu düşünmek için \emph{henüz} herhangi bir gerekçe
vermedik.}
\begin{align*}
	f(n+1) &=
	\begin{cases}
		a_n &\text{eğer }(\forall h \in S)\equivrep{h}{} \leq (a_n)_\Real\\
		f(n)&\text{aksi hâlde}
	\end{cases}\\
	g(n+1) &=
	\begin{cases}
		a_n &\text{eğer }(\exists h \in S)\equivrep{h}{} \geq (a_n)_\Real\\
	 	g(n) &\text{aksi hâlde}
	\end{cases}
\end{align*}
Hem $f$ hem $g$ birer Cauchy dizisidir. (Bunu denetlemek oldukça kolaydır,
fakat alıştırma olarak bırakıyoruz.) $f$ ile $g$ arasındaki fark her adımda
yarıya indiği için $(f-g)$ fonksiyonunun $0$'a yakınsadığına dikkat edin.
Dolayısıyla $\equivrep{f}{} = \equivrep{g}{}$ olur.

Önce $\equivrep{f}{}$'nin $S$ için bir üst sınır olduğunu, yani
$(\forall h \in S)\equivrep{h}{} \leq \equivrep{f}{}$ olduğunu gösterelim.
(İlerlerken \olref{thm:cauchyorderedfield}'e başvuracağız.) $h \in S$ olsun ve
olmayana ergiyle $\equivrep{f}{} < \equivrep{h}{}$ olduğunu varsayalım;
öyleyse $0_\Real < \equivrep{(h-f)}{}$ olur. $f$ monoton azalan bir Cauchy
dizisi olduğundan
$\equivrep{(c_{f(n)} - f)}{} < \equivrep{(h-f)}{}$ olacak biçimde bir
$n \in \Nat$ vardır. Böylece:
\[
	(f(n))_\Real = \equivrep{c_{f(n)}}{} < \equivrep{f}{} + \equivrep{(h-f)}{} = \equivrep{h}{},
\]
oysa kuruluş gereği $\equivrep{h}{} \leq (f(n))_\Real$'dir; çelişki.

Şimdi $\equivrep{f}{} = \equivrep{g}{}$'nin $S$ üzerindeki \emph{en küçük}
üst sınır olduğunu gösterelim. $j$ herhangi bir Cauchy dizisi olsun ve
$\equivrep{j}{} < \equivrep{g}{}$ olduğunu varsayın. Yukarıdakine benzer
biçimde ($g$'nin \emph{artan} olduğunu kullanarak)
$\equivrep{j}{} < (g(n))_\Real$ olacak biçimde bir $n \in \Nat$ vardır.
Ancak kuruluş gereği $(g(n))_\Real \leq \equivrep{h}{}$ olacak biçimde bir
$h \in S$ vardır; dolayısıyla
$\equivrep{j}{} < \equivrep{h}{}$ ve böylece $\equivrep{j}{}$, $S$ için bir
üst sınır değildir; burada üst sınırlar yine bu dizilerin denklik sınıfları için
anlaşılmaktadır.
\end{proof}

\end{document}
